% PRG v7: Forecast-timing conventions and the value of overnight information
% Target: Finance Research Letters (letter format: body <= 2,500 words, abstract <= 250 words)
% v7 rewrite 2026-07-14: headline = timing-convention flip (dual-convention framework)
% Evidence: K1699 (close, pinned), K1710 (mixed + open, pinned), K1544/K880 as directional anchors
% Authors: Yi-Hao Lai (Dayeh University)

\documentclass[11pt,a4paper]{article}

\usepackage[utf8]{inputenc}
\usepackage[T1]{fontenc}
\usepackage{mathptmx}           % Times Roman
\usepackage{amsmath,amssymb}
\usepackage{booktabs}
\usepackage{graphicx}
\usepackage{natbib}
\usepackage[margin=2.5cm]{geometry}
\usepackage{setspace}
\usepackage{hyperref}
\usepackage{caption}

\doublespacing
\hypersetup{colorlinks=true, citecolor=black, linkcolor=black, urlcolor=black, pdfborder={0 0 0}}

\title{Forecast-Timing Conventions and the Value \\ of Overnight Information in Volatility Forecasting\thanks{%
Computational infrastructure was provided by the author's research group.
All data and replication code are available upon request.}}

\author{Yi-Hao Lai\thanks{%
Department of Finance, Dayeh University, Changhua, Taiwan.
Email: yhlai@mail.dyu.edu.tw}}

\date{July 2026}

\begin{document}

\maketitle

\begin{abstract}
Session-level volatility models produce overnight and intraday forecast components that are naturally issued at different times: the overnight component at the previous close, the intraday component at the market open, after the overnight return has realized. We show that the forecast-timing convention---when each component is issued and what the benchmark observes---drives the apparent value of session-level modeling. Using a parsimonious Periodic Realized GARCH (PRG, 6--8 parameters) on pinned data for six equity, commodity, and futures markets, we evaluate the same model under three conventions. Under a mixed-timing design that session-level evaluation easily slips into---model components issued at two different times, benchmark held at the previous close---PRG beats GJR-GARCH with Diebold--Mariano statistics of $+4.3$ to $+6.4$ in all six markets. Under a coherent close-time convention, with every component issued at the previous close, the advantage vanishes: zero of six markets significant. Under a coherent open-time convention, where an information-matched GJR-X benchmark also observes the realized overnight return, PRG's advantage is positive in all six markets and on the pinned vintage exceeds the conservative $|t|>3$ threshold in five; the exception has the smallest overnight variance share. The same model and data thus deliver DM statistics from $-2.3$ to $+10.1$ depending only on the evaluation clock. Overnight information has genuine forecasting value at the open horizon and none at the day-ahead horizon; mixed-timing comparisons attribute to the model what belongs to the clock.

\medskip\noindent
\textbf{Keywords:} Volatility forecasting; Forecast timing; Overnight information; Periodic GARCH; Forecast evaluation

\medskip\noindent
\textbf{JEL Classification:} C22, C52, C53, G17
\end{abstract}

\newpage

%% =====================================================================
%% 1. INTRODUCTION
%% =====================================================================
\section{Introduction}

Financial markets alternate between an overnight period, when exchanges are closed yet information arrives continuously, and an intraday period of active price discovery. \citet{Blanc2014} document that the two sessions exhibit fundamentally different volatility feedback structures, and a substantial literature exploits session-level information for forecasting: periodic GARCH structures \citep{Bollerslev1996}, coupled intraday--overnight score-driven systems \citep{Linton2020}, continuous-time overnight diffusions \citep{Kim2023}, realized-variance models augmented with overnight returns \citep{Todorova2014,Opschoor2021}, and regime-switching session models \citep{Lai2024}. These studies generally report that exploiting the session split improves volatility forecasts.

This paper examines a degree of freedom in such comparisons that has received little systematic attention: \emph{when} each forecast component is issued. A session-level model naturally produces its overnight forecast at the previous close and its intraday forecast at the market open---after the overnight return has realized. Its full-day forecast is therefore a composite of components measurable with respect to two different information sets. Benchmarks such as GJR-GARCH, by contrast, issue a single day-ahead forecast at the previous close. Comparing the two-time composite against the one-time benchmark mixes a modeling question (does session structure help?) with an information question (what does acting at the open buy?). We show the mixture is not innocuous: it can manufacture the appearance of decisive model superiority where, at matched information sets, there is none.

Our vehicle is the Periodic Realized GARCH (PRG), a deliberately parsimonious model that extends the periodic GARCH of \citet{Bollerslev1996} from calendar to session frequency: a single variance recursion alternates between session-specific parameter sets (6--8 parameters), so each session's variance carries into the next. On pinned data for six markets---SPY, QQQ, GLD, EEM, the Taiwan 0050 ETF, and TAIFEX index futures---we evaluate the identical model under three forecast-timing conventions. Under the \emph{mixed} convention (model components at two times, benchmark at the close), PRG appears to dominate GJR-GARCH, with Diebold--Mariano statistics of $+4.3$ to $+6.4$ across all six markets. Under a coherent \emph{close-time} convention, all components issued at the previous close, the advantage disappears in all six markets. Under a coherent \emph{open-time} convention, against an information-matched GJR-X benchmark that also observes the realized overnight return, PRG wins in all six markets---five above the conservative $|t|>3$ threshold on the pinned vintage, with the direction stable across vintages.

The contribution is threefold. First, we quantify, across six markets on a single pinned data vintage, how far the forecast-timing convention alone moves headline test statistics---from $-2.3$ to $+10.1$.
Second, we provide a coherent dual-convention evaluation template: a close-time comparison for day-ahead applications and an information-matched open-time comparison for real-time updating. Third, the substantive finding is an honest localization of where overnight information pays: session-level modeling adds statistically significant value when forecasts are updated at the open---subject to a functional-form caveat we state explicitly in Section~\ref{sec:results}---and adds nothing over a plain GJR-GARCH for day-ahead forecasting.

%% =====================================================================
%% 2. MODEL AND FORECAST-TIMING CONVENTIONS
%% =====================================================================
\section{Model and forecast-timing conventions}

\subsection{Session decomposition and the PRG model}

Each trading day $d$ splits into an overnight session ($s=0$) and an intraday session ($s=1$). With $P^{c}_d$ and $P^{o}_d$ the close and open prices,
\begin{equation}
    r_{d,0} = \log(P^{o}_d / P^{c}_{d-1}), \qquad
    r_{d,1} = \log(P^{c}_d / P^{o}_d),
    \label{eq:sessions}
\end{equation}
and the common evaluation target for all models is the full-day variance proxy $\sigma^2_{\text{full},d} = r^2_{d,0} + r^2_{d,1}$. Under QLIKE loss, model rankings are robust to noise in unbiased proxies \citep{Patton2011,Hansen2006}. For TAIFEX, tick-level 5-minute realized variances replace squared session returns.

The PRG operates a single variance recursion over sessions indexed $n = 1, \ldots, 2D$ with session indicator $s_n$:
\begin{equation}
    h_n = \omega_{s_n} + \alpha_{s_n} x_{n-1} + \gamma_{s_n} x_{n-1}\,\mathbb{1}(r_{n-1}<0) + \beta_{s_n} h_{n-1},
    \label{eq:prg}
\end{equation}
where $x_n = r^2_n$ (the Basic variant sets $\gamma_s = 0$; ``Realized'' denotes the observed session-frequency proxy entering the recursion, distinct from the daily Realized GARCH of \citealp{Hansen2012}). Parameters are estimated by Gaussian QML \citep{BollerslevWooldridge1992}. With $\rho_s = \alpha_s + 0.5\gamma_s + \beta_s$, covariance stationarity requires only $\rho_0\rho_1 < 1$ \citep{Bollerslev1996}. The cross-session term $\beta_{s_n} h_{n-1}$ is the \emph{session bridge}: each session's conditional variance feeds the next.

\subsection{Three forecast-timing conventions}
\label{sec:conventions}

Let $\mathcal{F}^{c}_{d-1}$ be the information set at the day-$(d-1)$ close and $\mathcal{F}^{o}_{d}$ the set at the day-$d$ open; $\mathcal{F}^{o}_{d}$ contains the realized overnight return $r_{d,0}$. The PRG session forecasts are
\begin{align}
    \hat{h}_{d,0} &= \mathbb{E}\!\left[r^2_{d,0} \mid \mathcal{F}^{c}_{d-1}\right],
    \label{eq:h_ov} \\
    \hat{h}_{d,1}(z) &= \omega_{1} + \alpha_{1} z + \gamma_{1} z\,\mathbb{1}(r_{d,0}<0) + \beta_{1} \hat{h}_{d,0},
    \label{eq:h_in}
\end{align}
where $z$ is whatever stands in for the overnight squared return. Three full-day forecast objects follow:

\begin{itemize}
\item \textbf{Close convention (coherent).} Everything is issued at the day-$(d-1)$ close: $\hat{\sigma}^2_{\text{C}} = \hat{h}_{d,0} + \hat{h}_{d,1}(\hat{h}_{d,0})$, with the unobserved overnight square replaced by its model-consistent expectation (leverage indicator at its symmetric expectation $\tfrac12$). Every input is $\mathcal{F}^{c}_{d-1}$-measurable. The natural benchmark is GJR-GARCH \citep{Glosten1993} on close-to-close returns, also issued at $\mathcal{F}^{c}_{d-1}$.
\item \textbf{Open convention (coherent).} The full-day object is formed at the day-$d$ open: $\hat{\sigma}^2_{\text{O}} = r^2_{d,0} + \hat{h}_{d,1}(r^2_{d,0})$, with the realized overnight component plugged in. An information-\emph{matched} benchmark must also observe $r_{d,0}$: we use GJR-X, a daily GJR augmented with the \emph{current-day} overnight squared return,
$h_t = \omega + \alpha r^2_{t-1} + \gamma\,\mathbb{1}(r_{t-1}<0)\,r^2_{t-1} + \beta h_{t-1} + \delta\, r^2_{\text{overnight},t}$,
issued at the same open.
\item \textbf{Mixed convention (incoherent).} The composite $\hat{\sigma}^2_{\text{M}} = \hat{h}_{d,0} + \hat{h}_{d,1}(r^2_{d,0})$ combines an $\mathcal{F}^{c}_{d-1}$-measurable component with an $\mathcal{F}^{o}_{d}$-measurable one, and is compared against a benchmark held at $\mathcal{F}^{c}_{d-1}$. The composite corresponds to no coherent single-issuance-time forecast---it treats $r_{d,0}$ as unobserved in its first component yet observed in its second, choices no single information set makes simultaneously rational---and it is benchmarked against a model held at the earlier set $\mathcal{F}^{c}_{d-1}$, which is handicapped by construction. This is the convention implicit in evaluations that sum session forecasts ``as produced'' and compare against day-ahead models---including earlier drafts of this paper.
\end{itemize}

The mixed convention is not a straw man: session-level models naturally produce their components at different issuance times, single-recursion benchmarks are issued once a day, and nothing in standard evaluation practice forces the two information sets to be matched---earlier drafts of this paper reported exactly this comparison. Our point is that the resulting comparison confounds model structure with information timing, and Section~\ref{sec:results} quantifies the confound.

\subsection{Evaluation}
\label{sec:eval}

Losses are QLIKE; pairwise comparisons use \citet{Diebold1995} tests with Newey--West HAC variance (bandwidth $\lceil n^{1/3}\rceil$) and the \citet{Harvey1997} small-sample correction. With the 18 pairwise DM tests of Table~\ref{tab:flip} as the reported family ($\alpha/m = 0.05/18 \approx 0.0028$, two-sided $|z| \approx 2.99$) we adopt $|t|>3.0$ as a conservative, Bonferroni-style significance threshold; the six lagged-variant robustness tests discussed in Section~\ref{sec:results} are far from significance under any threshold, and we report nominal $p$-values throughout so readers can apply their own standard. All estimation is strictly out-of-sample with periodic refitting (GJR and GJR-X every 63 trading days; PRG every 126--252 sessions), using only data before each forecast origin.

%% =====================================================================
%% 3. DATA
%% =====================================================================
\section{Data}

Table~\ref{tab:data} summarizes the six markets. All OHLC series are dividend- and split-adjusted prices pinned as snapshot files (vintage 2026-07-12), so overnight returns exclude predictable ex-dividend jumps; every number in this paper reproduces bit-identically from the archived snapshots.\footnote{Earlier circulated drafts reported results from unpinned, vintage-drifting downloads (e.g., an SPY headline DM of 6.00 under the mixed convention). All results here derive from the pinned snapshots; the replication package documents the vintage, adjustment convention, and estimation seeds.}

\begin{table}[htbp]
\centering
\caption{Data and session decomposition}
\label{tab:data}
\small
\begin{tabular}{llrlc}
\toprule
Market & Source & OOS obs & OOS period & ON var share (\%) \\
\midrule
% source: experiments/k1699/README.md §3.4 pinned snapshot table (N, OOS periods)
% source (ON share): experiments/K1710/K1710_results.json .markets.<M>.oos_overnight_variance_share (×100, 1dp)
SPY        & OHLC (pinned)   & 1{,}823 & 2019-01--2026-04 & 44.8 \\
QQQ        & OHLC (pinned)   & 1{,}981 & 2018-05--2026-04 & 38.5 \\
GLD        & OHLC (pinned)   & 1{,}613 & 2019-10--2026-04 & 60.9 \\
EEM        & OHLC (pinned)   & 1{,}734 & 2019-05--2026-04 & 70.7 \\
0050.TW    & OHLC (pinned)   & 1{,}251 & 2021-01--2026-04 & 63.5 \\
TAIFEX TX  & Tick (5-min RV) & 843     & 2022-07--2025-12 & 68.9 \\
\bottomrule
\end{tabular}
\begin{flushleft}
\footnotesize
\emph{Notes:} OOS obs is the out-of-sample day count (70/30 split for OHLC markets). ON var share is the overnight share of full-day variance over the OOS window, computed from the pinned snapshots. TAIFEX uses tick-level 5-minute realized variance for the overnight (15:00--08:45) and regular (08:45--13:45) sessions of the volume-selected near-month contract; 0050.TW is cleaned for the 2014 split.
\end{flushleft}
\end{table}

%% =====================================================================
%% 4. RESULTS
%% =====================================================================
\section{Results}
\label{sec:results}

\subsection{The flip}

Table~\ref{tab:flip} is the paper's central result: the same PRG specification, on the same pinned data, evaluated under the three conventions of Section~\ref{sec:conventions}.

\begin{table}[htbp]
\centering
\caption{The forecast-timing flip: DM statistics for the same model and data under three conventions}
\label{tab:flip}
\small
\begin{tabular}{lrccc}
\toprule
 & & Mixed & Close & Open \\
 & & PRG vs.\ GJR & PRG vs.\ GJR & PRG vs.\ GJR-X \\
Market & $N$ & (incoherent) & (both at $\mathcal{F}^{c}_{d-1}$) & (both at $\mathcal{F}^{o}_{d}$) \\
\midrule
% Generated by scripts/gen_flip_table.py — do not hand-edit numbers.
% source (close col): experiments/k1699/k1699_results.json .markets.<M>.dm_tests.PRG_tminus1_exp_vs_GJR
% source (mixed col): experiments/K1710/K1710_results.json .markets.<M>.dm_tests.mixed_anchor_main
% source (open col):  experiments/K1710/K1710_results.json .markets.<M>.dm_tests.open_panel_main
%   JSON orientation in both files: negative t = PRG better; paper reports positive = PRG better (sign flipped)
SPY     & 1,823 & $+5.83^{***}$ ($<$0.001) & $-0.74$ (0.46) & $+3.56^{***}$ ($<$0.001) \\
QQQ     & 1,981 & $+4.78^{***}$ ($<$0.001) & $-2.28$ (0.02) & $+1.56$ (0.12) \\
GLD     & 1,613 & $+6.11^{***}$ ($<$0.001) & $+0.44$ (0.66) & $+3.64^{***}$ ($<$0.001) \\
EEM     & 1,734 & $+6.40^{***}$ ($<$0.001) & $+0.54$ (0.59) & $+10.14^{***}$ ($<$0.001) \\
0050.TW & 1,251 & $+5.19^{***}$ ($<$0.001) & $+0.32$ (0.75) & $+3.67^{***}$ ($<$0.001) \\
TAIFEX  & 843 & $+4.33^{***}$ ($<$0.001) & $+0.49$ (0.62) & $+5.50^{***}$ ($<$0.001) \\
\midrule
Harvey-significant ($|t|>3$) & & 6/6 & 0/6 & 5/6 \\
\bottomrule
\end{tabular}
\begin{flushleft}
\footnotesize
\emph{Notes:} PRG Extended specification (Eq.~\ref{eq:prg}) throughout. DM $t$-statistics on QLIKE loss differentials, reported as benchmark loss minus PRG loss, so positive values favor PRG; nominal two-sided $p$-values in parentheses; $^{***}$ marks the conservative $|t|>3.0$ threshold. Mixed: PRG composite $\hat{h}_{d,0} + \hat{h}_{d,1}(r^2_{d,0})$ vs.\ GJR issued at the close. Close: PRG composite $\hat{h}_{d,0} + \hat{h}_{d,1}(\hat{h}_{d,0})$ vs.\ GJR, both issued at the day-$(d-1)$ close. Open: PRG composite $r^2_{d,0} + \hat{h}_{d,1}(r^2_{d,0})$ vs.\ GJR-X with the current-day overnight regressor, both issued at the day-$d$ open. DM variances use Newey--West HAC with bandwidth $\lceil n^{1/3}\rceil$; doubling the bandwidth changes no verdict. All three panels use the identical pinned data vintage and deterministic estimation seeds.
\end{flushleft}
\end{table}

Read across any row. Under the mixed convention, PRG appears to dominate GJR decisively in every market---exactly the kind of headline the overnight-information literature reports. Under the coherent close-time convention, the advantage evaporates: no market is significant even at nominal levels, DM statistics straddle zero, and the only market approaching conventional significance (QQQ, $t=-2.28$) points \emph{against} PRG. Under the coherent open-time convention---where the benchmark is granted the same realized overnight return---PRG's advantage is positive in all six markets and exceeds $|t|>3$ in five; the exception is QQQ ($t=+1.56$), the market with the smallest overnight variance share. The spread between conventions within a single market's row reaches $9.6$ $t$-units (EEM). Nothing about the model or the data changes across the panels; only the clock does.

\subsection{Interpretation}

\paragraph{What mixed-timing comparisons measure.} The mixed panel's inflation has a mechanical source: the intraday component $\hat{h}_{d,1}(r^2_{d,0})$ absorbs the realized overnight shock while the GJR benchmark, stuck at the previous close, cannot. The gap between the mixed and close panels---$3.8$ to $7.1$ $t$-units per market---is the part of the headline that belongs to the information clock, not to session modeling. Studies that sum session-level forecasts ``as issued'' and benchmark them against day-ahead models inherit this confound, and their reported significance levels should be read accordingly.

\paragraph{Where overnight information genuinely pays.} The open panel shows that the value of session modeling at the open horizon is real, not an artifact of information mismatch: GJR-X sees the same realized overnight return, yet PRG's session-specific propagation of that information---through parameters estimated separately for the overnight and intraday processes and linked by the bridge---beats the exogenous-regressor route in all six markets. The cross-market pattern is itself informative: on the pinned vintage the open-time DM statistic is perfectly rank-ordered in the market's overnight variance share, from EEM (70.7\%, $t=+10.14$) down to QQQ (38.5\%, $t=+1.56$, the only insignificant market)---descriptive with six markets, vintage-sensitive at the low-share end (see Robustness below), but directionally the ordering a genuine overnight-information channel predicts. One caveat keeps this honest: part of the gap may reflect functional form rather than information processing, since PRG embeds the realized overnight component additively in its composite while GJR-X must load it through a single linear regressor. An intraday-target-only design that removes the additive component from both sides is left for follow-up.

\paragraph{No exception in high-overnight-share markets.} A natural conjecture is that markets whose overnight session carries most of the variance (TAIFEX at 68.9\%, EEM at 70.7\%, GLD at 60.9\%) might retain a close-time PRG advantage. They do not: their close-time DM statistics are $+0.49$, $+0.54$, and $+0.44$---directionally favorable, nowhere near significance.
% source: experiments/k1699/k1699_results.json close-panel t_stats (sign flipped), README §4.3 point 2

\paragraph{Robustness.} The close-time null is robust to how the unobserved overnight square is plugged in: a lagged-realized variant ($z = r^2_{d-1,0}$) reproduces every sign and every insignificance of the expectation variant.
% source: experiments/k1699/k1699_results.json .markets.<M>.dm_tests.PRG_tminus1_lag_vs_GJR (0/6 Harvey)
Doubling the HAC bandwidth changes no verdict in any panel. Point statistics---QLIKE levels and individual DM values---are sensitive to data vintage and to estimation basins of the non-convex likelihood: an earlier unpinned run of the open-time comparison put SPY at $+2.1$ and QQQ at $+3.0$ versus the pinned $+3.56$ and $+1.56$ here. No direction flips in any market under any vintage, and the pinned values reproduce bit-identically from the archived snapshots. Two headline framings are themselves vintage-sensitive and should be read as pinned-vintage statements: the open-panel Harvey count (five of six here; four of six in the pilot vintage, where SPY sits below the threshold) and the perfect share-ordering (the pilot inverts the two lowest-share markets). The vintage-robust statements are the direction---positive in all six markets---and the large-margin, high-share markets (EEM, TAIFEX); we make no claims from raw QLIKE rankings or borderline point values.

%% =====================================================================
%% 5. DISCUSSION AND CONCLUSION
%% =====================================================================
\section{Discussion and conclusion}

This paper's finding is a methodological warning with a constructive core. The warning: forecast-timing conventions are not a technical footnote. The same session-level volatility model, on the same pinned data, delivers Diebold--Mariano statistics anywhere from $-2.3$ to $+10.1$ depending only on when forecast components are issued and what the benchmark observes. Comparisons in the overnight-information literature that combine components issued at different times---or that benchmark open-informed forecasts against close-informed models---can overstate model value by several $t$-units per market. Reported significance under such designs conflates the model with the clock. The constructive core: evaluated coherently, overnight information has a well-defined and horizon-specific value. For day-ahead forecasting at the previous close, session-level structure adds nothing over a plain GJR-GARCH: the overnight return is not yet observed, and forecasting it session-by-session is no better than forecasting the day directly. For real-time updating at the open, session structure adds value that an exogenous-regressor patch does not replicate---significantly so in five of six markets on the pinned vintage, increasingly so the larger the market's overnight variance share, and subject to the functional-form caveat of Section~\ref{sec:results}.

The practical prescription for researchers is to report the issuance time of every forecast component and match the benchmark's information set to it---our close/open panel pair provides a template. The prescription for practitioners is equally direct: a session-level model such as PRG is a tool for the open, not for the close. PRG itself remains attractive as a vehicle---6--8 parameters, standard QML, no latent regimes \citep{Lai2024,Haas2004}---precisely because its parsimony makes the convention experiment clean: nothing in the flip can be attributed to estimation complexity.

Limitations: OHLC markets use squared session returns as variance proxies (QLIKE rankings are proxy-robust, \citealp{Patton2011}); tail-risk and economic-value evaluation under the two coherent conventions, and an intraday-target-only comparison isolating functional form from information, are natural follow-ups.

%% =====================================================================
%% REFERENCES
%% =====================================================================
\bibliographystyle{apalike}

\begin{thebibliography}{99}

\bibitem[Blanc et~al.(2014)]{Blanc2014}
Blanc, P., Chicheportiche, R., and Bouchaud, J.-P. (2014).
\newblock The fine structure of volatility feedback II: Overnight and intra-day effects.
\newblock \emph{Physica A}, 402, 58--75.
\newblock \url{https://doi.org/10.1016/j.physa.2014.01.047}

\bibitem[Bollerslev and Ghysels(1996)]{Bollerslev1996}
Bollerslev, T. and Ghysels, E. (1996).
\newblock Periodic autoregressive conditional heteroskedasticity.
\newblock \emph{Journal of Business \& Economic Statistics}, 14(2), 139--151.
\newblock \url{https://doi.org/10.1080/07350015.1996.10524640}

\bibitem[Bollerslev and Wooldridge(1992)]{BollerslevWooldridge1992}
Bollerslev, T. and Wooldridge, J.~M. (1992).
\newblock Quasi-maximum likelihood estimation and inference in dynamic models with time-varying covariances.
\newblock \emph{Econometric Reviews}, 11(2), 143--172.
\newblock \url{https://doi.org/10.1080/07474939208800229}

\bibitem[Diebold and Mariano(1995)]{Diebold1995}
Diebold, F.~X. and Mariano, R.~S. (1995).
\newblock Comparing predictive accuracy.
\newblock \emph{Journal of Business \& Economic Statistics}, 13(3), 253--263.
\newblock \url{https://doi.org/10.1080/07350015.1995.10524599}

\bibitem[Glosten et~al.(1993)]{Glosten1993}
Glosten, L.~R., Jagannathan, R., and Runkle, D.~E. (1993).
\newblock On the relation between the expected value and the volatility of the nominal excess return on stocks.
\newblock \emph{The Journal of Finance}, 48(5), 1779--1801.
\newblock \url{https://doi.org/10.1111/j.1540-6261.1993.tb05128.x}

\bibitem[Haas et~al.(2004)]{Haas2004}
Haas, M., Mittnik, S., and Paolella, M.~S. (2004).
\newblock A new approach to Markov-switching GARCH models.
\newblock \emph{Journal of Financial Econometrics}, 2(4), 493--530.
\newblock \url{https://doi.org/10.1093/jjfinec/nbh020}

\bibitem[Hansen and Lunde(2006)]{Hansen2006}
Hansen, P.~R. and Lunde, A. (2006).
\newblock Consistent ranking of volatility models.
\newblock \emph{Journal of Econometrics}, 131(1--2), 97--121.
\newblock \url{https://doi.org/10.1016/j.jeconom.2005.01.005}

\bibitem[Hansen et~al.(2012)]{Hansen2012}
Hansen, P.~R., Huang, Z., and Shek, H.~H. (2012).
\newblock Realized GARCH: A joint model for returns and realized measures of volatility.
\newblock \emph{Journal of Applied Econometrics}, 27(6), 877--906.
\newblock \url{https://doi.org/10.1002/jae.1234}

\bibitem[Harvey et~al.(1997)]{Harvey1997}
Harvey, D., Leybourne, S., and Newbold, P. (1997).
\newblock Testing the equality of prediction mean squared errors.
\newblock \emph{International Journal of Forecasting}, 13(2), 281--291.
\newblock \url{https://doi.org/10.1016/S0169-2070(96)00719-4}

\bibitem[Kim et~al.(2023)]{Kim2023}
Kim, D., Shin, M., and Wang, Y. (2023).
\newblock Overnight GARCH-It\^{o} volatility models.
\newblock \emph{Journal of Business \& Economic Statistics}, 41(4), 1215--1227.
\newblock \url{https://doi.org/10.1080/07350015.2022.2116027}

\bibitem[Lai et~al.(2024)]{Lai2024}
Lai, Y.-H., Wang, Y.-C., and Chang, Y.-C. (2024).
\newblock Forecasting trading-session return volatility in Taiwan futures market:
A periodic regime switching with jump approach.
\newblock \emph{Asia-Pacific Financial Markets}, 31(2), 285--305.
\newblock \url{https://doi.org/10.1007/s10690-023-09415-w}

\bibitem[Linton and Wu(2020)]{Linton2020}
Linton, O. and Wu, J. (2020).
\newblock A coupled component DCS-EGARCH model for intraday and overnight volatility.
\newblock \emph{Journal of Econometrics}, 217(1), 176--201.
\newblock \url{https://doi.org/10.1016/j.jeconom.2019.12.015}

\bibitem[Opschoor and Lucas(2021)]{Opschoor2021}
Opschoor, A. and Lucas, A. (2021).
\newblock Observation-driven models for realized variances and overnight returns applied to Value-at-Risk and Expected Shortfall forecasting.
\newblock \emph{International Journal of Forecasting}, 37(2), 622--633.
\newblock \url{https://doi.org/10.1016/j.ijforecast.2020.07.009}

\bibitem[Patton(2011)]{Patton2011}
Patton, A.~J. (2011).
\newblock Volatility forecast comparison using imperfect volatility proxies.
\newblock \emph{Journal of Econometrics}, 160(1), 246--256.
\newblock \url{https://doi.org/10.1016/j.jeconom.2010.03.034}

\bibitem[Todorova and Sou\v{c}ek(2014)]{Todorova2014}
Todorova, N. and Sou\v{c}ek, M. (2014).
\newblock Overnight information flow and realized volatility forecasting.
\newblock \emph{Finance Research Letters}, 11(4), 420--428.
\newblock \url{https://doi.org/10.1016/j.frl.2014.07.001}

\end{thebibliography}

\end{document}
